% STATUS: written 2026-11-05
% Sources read in full: Raudenbush 1934 (Trans. AMS 36, 361--368,
% Theorems 1, 4--8); Wu 1987/1989 (Bull. AMS 20, 175--184, zero
% decomposition); Hubbard-Lundell 2011 (Amer. Math. Monthly 118,
% 245--261, sec. 4).
\section{The Ritt--Kolchin theory}\label{sec:differential-algebra:ritt-kolchin}

Chapter~\ref{ch:origins} defined the differential field and the
differential ideal with Raudenbush
(Definition~\ref{def:origins:differential-field}). Ritt asked how far the
ideal theory of these rings goes \cite[p.~361]{Raudenbush1934}.
Raudenbush rebuilt the theory in 1934 to give the theorem of Ritt a
clean proof \cite[p.~361]{Raudenbush1934}. He replaced the differential
ideal by its radical inside
the differential category: the perfect differential ideal
\cite[p.~361]{Raudenbush1934}.

A \emph{perfect} differential ideal is one that contains an element
whenever it contains a positive power of it
\cite[p.~362]{Raudenbush1934}. In the algebraic theory the perfect
differential ideals are the radicals of differential ideals
\cite[p.~362]{Raudenbush1934}. Raudenbush first shows that the perfect
ideal determined by a set is the radical of the differential ideal it
determines.

\begin{theorem}[{\cite[Theorem~1, p.~363]{Raudenbush1934}}]\label{thm:diffalg:perfect-is-radical}
Let $R$ be a differential ring of characteristic zero and let
$\alpha$ be a set of its elements. If $t$ is any element of the perfect
differential ideal $\{\alpha\}$ determined by $\alpha$, then some
positive integral power of $t$ is in the differential ideal $[\alpha]$
determined by $\alpha$.
\end{theorem}

The proof uses the product rule
\cite[p.~363]{Raudenbush1934}. From $ab$ in the ideal, Raudenbush first
shows every product $a_pb_q$ of a derivative of $a$ by a derivative of
$b$ is in the ideal \cite[p.~363]{Raudenbush1934}. Multiplying the
linear expressions of the two powers of $t$ against each other then
places each term in the smaller ideal \cite[p.~363]{Raudenbush1934}.
The same device gives the identity
$\{\alpha, a\}\cap\{\alpha,b\} = \{\alpha, ab\}$ for any two elements of
$R$ \cite[Theorem~2, p.~363]{Raudenbush1934}, and this identity drives
the decomposition theory \cite[p.~365]{Raudenbush1934}.

A perfect differential ideal $\pi$ is \emph{reducible} if it is the
intersection of two perfect differential ideals that both contain it
properly \cite[p.~365]{Raudenbush1934}. Otherwise it is
\emph{irreducible} \cite[p.~365]{Raudenbush1934}. The relations between
primality and irreducibility are the first structural facts (Theorem~\ref{thm:diffalg:perfect-decomposition})
\cite[p.~365]{Raudenbush1934}.\begin{theorem}[{\cite[Theorems~4 and~5, p.~365]{Raudenbush1934}}]\label{thm:diffalg:perfect-decomposition}
Let $R$ be a differential ring with a basis theorem: every perfect
differential ideal has a finite basis.

\begin{enumerate}
\item \emph{Irreducible implies prime}. A perfect differential ideal
  which is irreducible is prime.
\item \emph{Decomposition}. Any perfect differential ideal is the
  intersection of a finite set of irreducible, hence prime, perfect
  differential ideals.
\end{enumerate}
\end{theorem}

The decomposition theorem follows by noetherian induction
\cite[p.~365]{Raudenbush1934}. Suppose the claim fails; the failing
ideal is reducible, one component still fails, and so on. An ascending
chain of perfect differential ideals appears, and the basis theorem
forbids it \cite[p.~365]{Raudenbush1934}. Uniqueness comes next: delete
in turn each member of the finite set that contains a different member,
and the remainder is determined by the ideal (Theorem~\ref{thm:diffalg:essential-uniqueness})
\cite[pp.~365--366]{Raudenbush1934}.\begin{theorem}[{\cite[Theorem~6, p.~365]{Raudenbush1934}}]\label{thm:diffalg:essential-uniqueness}
If a perfect differential ideal is the intersection of each of two
essential sets of prime perfect differential ideals, then the two sets
have the same number of members and coincide after a suitable
rearrangement.
\end{theorem}

All this applies to the differential ring that motivates the
whole subject: the ring of forms (differential polynomials) in a finite
set of indeterminates over a differential field
\cite[p.~361]{Raudenbush1934}. Ritt had proved his basis theorem for
this ring \cite[p.~361]{Raudenbush1934}. Raudenbush gives a second,
shorter proof and remarks that its simplicity reflects the strength of
the theory (Theorem~\ref{thm:diffalg:ritt-raudenbush-basis}) \cite[p.~362]{Raudenbush1934}.

\begin{theorem}[{\cite[Theorem~7, pp.~366--367]{Raudenbush1934}}]\label{thm:diffalg:ritt-raudenbush-basis}
Let $F$ be a differential field of characteristic zero. Then the
differential ring of forms in a finite set of indeterminates over $F$ is
a differential ring with a basis theorem.
\end{theorem}

With the basis theorem in hand, the preceding theorems give the
intersection theorem on primes
\cite[Theorems~5, 6 and~7, pp.~365--367]{Raudenbush1934}. Every perfect
differential ideal of the ring of differential polynomials over a
characteristic-zero differential field is a finite intersection of prime
perfect differential ideals, unique up to rearrangement
\cite[Theorems~5, 6 and~7, pp.~365--367]{Raudenbush1934}. Combined with
Theorem~\ref{thm:diffalg:perfect-is-radical} this recovers the finite
basis theorem of Ritt for differential polynomials, and the finite decomposition
of differential algebraic varieties into irreducible components
\cite[p.~361]{Raudenbush1934}.

The method of Ritt himself is more explicit than the ideal theory,
because it
comes from elimination \cite[p.~176]{Wu1987RittConstructive}. For a finite system of differential polynomials he produces an
\emph{ascending set}, a finite ordered list of differential polynomials
of increasing class. He also produces a \emph{characteristic set} of the
differential ideal of the system, whose remainder theory decides membership
\cite[p.~176]{Wu1987RittConstructive}. Wu Wen-tsun rebuilt the whole
theory constructively around this device
\cite[pp.~175--184]{Wu1987RittConstructive}. His decomposition theorem
makes the irreducible components computable
\cite[p.~179]{Wu1987RittConstructive}. It is the computational form of the
prime decomposition of a perfect differential ideal stated above this
paragraph \cite[p.~179]{Wu1987RittConstructive}.

\begin{theorem}[{\cite[p.~179]{Wu1987RittConstructive}}]\label{thm:diffalg:wu-zero-decomposition}
Let $(\mathrm{DP})$ be a finite system of differential polynomials and
$G$ a differential polynomial. There is a mechanical procedure which
decides in a finite number of steps whether the zero set
$\mathrm{Zero}(\mathrm{DP}/G)$ is empty, and in the non-empty case
furnishes a decomposition
\[
\mathrm{Zero}(\mathrm{DP}/G) \;=\; \bigcup_i
\mathrm{Zero}(\mathrm{ASC}_i / I_i, G_i)
\]
where each $(\mathrm{ASC}_i)$ is a d-irreducible ascending set, $I_i$ is
the product of the initials of $(\mathrm{ASC}_i)$, and $G_i$ is reduced
with respect to $(\mathrm{ASC}_i)$.
\end{theorem}

The characteristic-set method is the computational side of the
Ritt--Kolchin theory \cite[pp.~175--184]{Wu1987RittConstructive}.
Modern textbook treatments \cite[pp.~245--261]{HubbardLundell2011}
develop exactly this interface between the abstract ideal theory and the
explicit elimination \cite[pp.~245--250]{HubbardLundell2011}.
Theorem~\ref{thm:diffalg:wu-zero-decomposition} is also what runs the
mechanical theorem proving in differential geometry that Wu proposed
\cite[pp.~175--184]{Wu1987RittConstructive}. The analogue on the Galois
side of the chapter is the algebraic group computation of
Section~\ref{sec:differential-algebra:differential-galois}.

% Moved from chapters/20-applications/03-differential-equations.tex (sources read: Kolchin 1973, Ax 1971 via Aslanyan 2018).
\section{Differential equations: fields with a derivation}\label{sec:differential-algebra:differential-equations}

The algebra of smooth functions on a manifold from
Section~\ref{sec:commutative:geometry} is a differential ring. Differential
algebra takes a differential equation as data and studies its solutions by
algebra, as Kolchin axiomatized it \cite{Kolchin1973}; Kolchin
began the theory in \cite[\S I.1]{Kolchin1973}.

\begin{definition}[{\cite[\S I.1, p.~58]{Kolchin1973}}]\label{def:diffalg:kolchin-diff-ring}
Let $\mathfrak{R}$ be a ring. An operator $\delta$ on $\mathfrak{R}$ is a
\emph{derivation operator} if $\delta(a+b)=\delta a+\delta b$ and
$\delta(ab)=(\delta a)b+a\,\delta b$ for all $a,b\in\mathfrak{R}$. A
\emph{differential ring} is a ring $\mathfrak{R}$ equipped with a finite set
$\Delta$ of derivation operators on $\mathfrak{R}$ such that
$\delta\delta'a=\delta'\delta a$ for every $a\in\mathfrak{R}$ and every
$\delta,\delta'\in\Delta$. If $\mathfrak{R}$ is a field, it is a
\emph{differential field}.
\end{definition}

Kolchin attaches to a differential ring the free commutative semigroup
$\Theta$ generated by $\Delta$ \cite[\S I.1, p.~59]{Kolchin1973}. An element
$\theta\in\Theta$ is a \emph{derivative operator}, and its \emph{order}
$\mathrm{ord}\,\theta$ is the sum of the exponents
\cite[\S I.1, p.~59]{Kolchin1973}. The objects of
differential algebra are then the differential subrings, differential ideals
and differential extensions, parallel to rings, ideals and field extensions
\cite[\S I.1, p.~59]{Kolchin1973}.

The differential equation itself becomes a statement about field extensions
\cite{Ax1971}.
Ax showed that the transcendence of a solution of the exponential differential
equation is bounded by the underlying derivation \cite[Theorem~3]{Ax1971}.
We recall his theorem from the account of Aslanyan \cite{Aslanyan2018}.

\begin{theorem}[{\cite[Theorem~1.1, p.~1]{Aslanyan2018}}]\label{thm:diffalg:ax-schanuel}
Let $K=(K;+,\cdot,D)$ be a differential field with field of constants $C$. If
$(x_1,y_1),\ldots,(x_n,y_n)$ are non-constant solutions to the exponential
differential equation $Dy=y\,Dx$, then
\[
\mathrm{td}_C\,C(\bar{x},\bar{y}) - \mathrm{l.\,dim}_{\mathbb{Q}}(\bar{x}/C) \geq 1,
\]
where $\mathrm{l.\,dim}_{\mathbb{Q}}(\bar{x}/C)$ is the dimension of the
$\mathbb{Q}$-span of $x_1,\ldots,x_n$ in the quotient vector space $K/C$.
\end{theorem}

This is the Ax--Schanuel theorem, the differential analogue of the open
conjecture of Schanuel in number theory; Ax \cite{Ax1971} proved it as an
application of the differential algebra of Kolchin \cite{Kolchin1973}. The
field of definition of a solution of a differential equation is read from the
bound \cite[\S2]{Aslanyan2018}. Section~\ref{sec:differential-algebra:transcendence}
comes back to it in the model-theoretic form the theorem of Ax started.

% Moved from chapters/20-applications/05-control.tex (source read in full: Fliess 1989).
\section{Control theory: differential fields and dynamic feedback}\label{sec:differential-algebra:control}

The engineer who designs an automatic controller writes a dynamical system as
$\dot x=f(x,u)$ with $u$ an input and $x$ a state \cite{Fliess1989}. Fliess observed in 1989 that
the input--output behaviour of the system lives in a finitely generated
extension of a differential field, in the sense of
Section~\ref{sec:differential-algebra:differential-equations} \cite{Fliess1989}.

\begin{definition}[{\cite[\S1.2, p.~228]{Fliess1989}}]\label{def:diffalg:fliess-systemic-field}
A \emph{systemic field} $\mathcal{K}$ is a finitely generated differential
extension of the base differential field $k$. There exists a finite tuple
$w=(w_1,\ldots,w_s)$ of elements of $\mathcal{K}$ with
$\mathcal{K}=k\langle w\rangle$, the differential field generated by $k$ and
$w$. An \emph{input} (or \emph{control}, \emph{command}) is any set containing
a differential transcendence basis of $\mathcal{K}/k$; if the set is itself
such a basis, the input is \emph{independent}.
\end{definition}

There is thus a partition $w=(u,y)$, not necessarily unique, into an
independent input $u$ and a remaining tuple $y$
\cite[\S1.2, p.~228]{Fliess1989}. The input $u$ is a
differential transcendence basis of $\mathcal{K}/k$, and the components of
$y$ are differentially algebraic over $k\langle u\rangle$
\cite[\S1.2, p.~228]{Fliess1989}. The state-space realisation question then
asks for a writing of the extension with a finite non-differential
transcendence basis \cite[\S1.2, p.~228]{Fliess1989}. Fliess shows that in general one cannot realise a
nonlinear system globally in explicit first-order form
\cite[Introduction, p.~227]{Fliess1989}.

\begin{definition}[{\cite[\S1.3, p.~228]{Fliess1989}}]\label{def:diffalg:io-system}
An \emph{input-output system} consists of an input $u=(w_1,\ldots,w_m)$ and an
output $y=(y_1,\ldots,y_p)$ with components in a universal differential field,
such that $y_1,\ldots,y_p$ are differentially algebraic over $k\langle u\rangle$.
If $w_1,\ldots,w_m$ are differential indeterminates, the input is
\emph{independent}.
\end{definition}

The framework applies to feedback synthesis, input-output inversion and
realization, and Fliess uses it to define a feedback as a specialisation of
differential fields \cite[Introduction, p.~227]{Fliess1989}.

\section{Automatic differentiation: derivations in computer code}\label{sec:differential-algebra:automatic-differentiation}

The topological refinement of Section~\ref{sec:cohomology:algebraic-topology}
computed derivations as homology of a chain complex. The engineering reading
of the same rule makes a computer program evaluate the derivative of the
function it implements \cite[Introduction, p.~2]{YuanThomsonAnderson2020}.
The method is automatic differentiation, and the program runs the chain rule
and the Leibniz rule step by step \cite{YuanThomsonAnderson2020}.
Yuan, Thomson and Anderson apply the method to the flight dynamics of a
tilt-rotor aircraft \cite{YuanThomsonAnderson2020}, where automatic
differentiation gives the stability derivatives.

The authors define automatic differentiation (AD) as a set of techniques for
evaluating numerically the derivative of a function implemented in a computer
program \cite[Introduction, p.~2]{YuanThomsonAnderson2020}.

\begin{definition}[{\cite[Introduction, p.~2]{YuanThomsonAnderson2020}}]\label{def:app:autodiff}
\emph{Automatic differentiation} (AD) is a set of techniques for numerically
evaluating the derivative of a function implemented within a computer
program. The method is based on the chain rule as applied to the
differentiation process.
\end{definition}

Every program decomposes into elementary operations
\cite[\S II.B]{YuanThomsonAnderson2020}. An arithmetic operator or
an intrinsic function applied to intermediate values is one elementary step
\cite[\S II.B]{YuanThomsonAnderson2020}. A step $g=h(a,b)$ of the program
differentiates by the chain rule as
\[
dg \;=\; \frac{\partial h}{\partial a}\,da \;+\; \frac{\partial h}{\partial b}\,db,
\]
\cite[equation (2)]{YuanThomsonAnderson2020}. Repeated use of the formula
simplifies the differentiation to a series of algebraic operations, with
derivative results coming from the elementary operations
\cite[\S II.B]{YuanThomsonAnderson2020}. In the forward mode of \cite[\S II.B]{YuanThomsonAnderson2020},
the propagation of the derivative of the state $\dot x_i$ with respect to an
initial condition $x_j$ is evaluated. The operation of the forward mode is
\[
\frac{\partial \dot x_i}{\partial x_j}
\;=\;
\sum_{k=1}^{n}\frac{\partial \dot x_i}{\partial x_k}\,
\frac{\partial x_k}{\partial x_j}
\;+\;
\sum_{l=1}^{m}\frac{\partial \dot x_i}{\partial u_l}\,
\frac{\partial u_l}{\partial x_j},
\]
with $n=\dim(x)$ and $m=\dim(u)$ \cite[equation
(3)]{YuanThomsonAnderson2020}.

The derivative of the composed function is thus read from the Jacobian
products, never from a difference quotient \cite[Introduction,
p.~2]{YuanThomsonAnderson2020}. Numerical differentiation by
finite differences would carry a truncation error and is avoided
\cite[Introduction, p.~2]{YuanThomsonAnderson2020}. The Leibniz rule of a
derivation on an algebra is one of the elementary operator identities that the
overloaded algebra of AD satisfies \cite[\S II.B]{YuanThomsonAnderson2020}.
The semantics of the overloaded arithmetic are those of the pair
$(x,\dot x)$, the state and its derivative, as in \cite[\S II.B]{YuanThomsonAnderson2020}. The composition is a chain rule
on the first-order jets of the elementary steps \cite[equations
(2)--(3)]{YuanThomsonAnderson2020}.

The axioms Messmer and Wood gave for the theory of separably closed fields with higher derivations are modified, with the modification stated through Witt vectors~\cite{WlWittVectorsAndSeparablyClosedF2019}.

The primitive element theorem is proved for differential fields whose derivation is not required to be constant on the base, which is the form the chapter's field theory needs~\cite{WlThePrimitiveElementTheoremForD2015}.

A nonlinear equation is solved in linear stages from its higher order derivations, which is the constructive use of a higher derivation the chapter records from the differential-algebra side~\cite{WlSolvingNonlinearEquationsFromH1990}.

For a differential field of characteristic zero with a basic set of $m$ derivations, the dimension polynomials of the intermediate extensions are determined~\cite{WlDimensionPolynomialsOfIntermed2020}.

Kondrat'eva tests a dimension polynomial of a differential field extension generated by a system of linear differential equations for minimality in the sense of Sit, and solves the problem in several particular cases~\cite{WlAMinimalDimensionPolynomialOfA1989}.

Let $K$ be a differential field over $\mathbb{C}$ with derivation $D$ and $G$ a finite linear automorphism group over $K$ preserving $D$. Tanabe shows that the finite-dimensional vertex algebra modules over the fixed point subfield $K^G$ are described in terms of the modules over $K$~\cite{WlFiniteDimensionalVertexAlgebra2012}.

Chalkley proves that the algebraic differential equation $P(w,w')=0$ of first order and second degree admits a degree-reducing transformation whenever a stated discriminant form has a representation over the differential field~\cite{WlAlgebraicDifferentialEquations1975}.

Marumo, Oaku and Takemura prove that the $n$-th power of a function satisfying a second-order linear differential equation satisfies an $(n+1)$-st order equation, give a method obtaining that equation, and determine its exponents and a bound on the degree of its polynomial coefficients~\cite{WlPropertiesOfPowersOfFunctionsS2015}.

Kondrat'eva studies a minimal dimension polynomial of the extension of a field given by a system of linear differential equations~\cite{WlAMinimalDimensionPolynomialOfT2020}.

Dimension polynomials of intermediate differential field extensions and the strength of the associated invariants are studied~\cite{WlDimensionPolynomialsOfIntermed2020}.

Bivariate dimension polynomials of non-reflexive prime difference-differential ideals are studied by characteristic sets~\cite{WlBivariateDimensionPolynomialsO2020}.
